\noindent A intercambialidade entre métodos de medição de material particulado depende da preservação de uma mesma base de comparação. Esta revisão sistematizada avaliou a comparabilidade entre métodos automáticos convencionais e gravimétricos para \MPdez{} e \MPdoiscinco{} no ar ambiente externo. Buscas em Scopus, Web of Science e PubMed, complementadas por rastreamento delimitado de citações, reuniram 437 registros únicos. Dos 84 relatos disponíveis para avaliação, 66 foram selecionados, abrangendo publicações de 1991 a 2025. A extração focal e a análise crítica descritiva distinguiram frações, configurações, comparadores, escalas temporais e definições estatísticas, sem metanálise. As diferenças não apresentaram direção fixa por tecnologia. O aquecimento de microbalanças frequentemente reduziu a massa medida, enquanto sistemas que compensaram perdas semivoláteis puderam superar filtros de referência sujeitos à volatilização. Atenuação beta e métodos ópticos mostraram dependência de umidade, propriedades do aerossol, concentração e tratamento dos dados. Correções melhoraram resultados em conjuntos específicos, com alcance condicionado ao desenho de validação. Correlação elevada não assegurou concordância, e razões de médias, médias de razões e incertezas não constituíram uma escala comum de desempenho. Os achados sustentam a avaliação do sistema completo de medição e a documentação das transições instrumentais antes de integrar séries ou avaliar representatividade temporal.
\vspace{.5em}
\noindent\textbf{Palavras-chave:} material particulado; gravimetria; monitoramento contínuo; comparabilidade; umidade relativa.
